\section{Formalization architecture}\label{sec:formalization}

The Lean~4 formalization is split between the shared repository package \texttt{lean-machine}
(\texttt{../lean-machine}) and this paper's package \texttt{BoundedArityDecoders} (\texttt{lean/}),
both pinned to Lean~4 and Mathlib version \texttt{v4.34.1} \cite{Mathlib2026}. Through Lake path
dependencies in \leanfile{lean/lakefile.toml}, \texttt{BoundedArityDecoders} requires
\leanname{SharedDecoders} (\cite{SharedDecoders2026}, which transitively imports
\leanname{SwapRigidity} from \cite{SeventeenthPaper2026} and \leanname{Transpose} from
\cite{FirstPaper2026}), \leanname{ClassUniformMemo} (\cite{ClassUniformMemo2026}), and
\leanname{LeanMachine}. None of the earlier papers is modified, and for this paper
\texttt{lean-machine} is extended only by registering four reusable modules under
\leanfile{LeanMachine/ShuffleShear/} in the package root, axiom audit, and \texttt{README}:
\begin{itemize}[leftmargin=1.5em]
\item \leanfile{Model.lean}, \leanfile{StoneShuffle.lean}, \leanfile{StoneUnshuffle.lean}, and
  \leanfile{RowShear.lean}: standalone deterministic multitape Turing machine semantics and verified
  linear-time machines for Stone's shuffle ($3$ work tapes, $\le 150 L$ steps), Stone's unshuffle
  ($2$ work tapes, $\le 150 L$ steps), and within-row cyclic shear ($2$ work tapes, $\le 200 L$
  steps).
\end{itemize}

Inside \leanfile{lean/BoundedArityDecoders/}, twenty-seven public modules are organized into three layers:
\begin{enumerate}[label=(\arabic*)]
\item \emph{Nine definition and specification modules}
  (\leanfile{AffineCubeComposition.lean},
  \leanfile{SumsetRankDeficit.lean},
  \leanfile{ExtensionLinearRun.lean},
  \leanfile{ExtensionAperiodicSix.lean},
  \leanfile{ExtensionTwoSidedProfile.lean},
  \leanfile{SwapCodeAnswers.lean},
  \leanfile{OperatorDuality.lean},
  \leanfile{NativeSweepInterface.lean}, and
  \leanfile{BoundedAritySetCapacity.lean}), defining the algebraic operators, one-stage linear
  run-extension invariants, coarse--fine interval decompositions, syntactic causal cones, and named
  proposition types;
\item \emph{Eleven primary proof modules}
  (\leanfile{AffineCubeCompositionProofs.lean},
  \leanfile{SumsetRankDeficitProofs.lean},
  \leanfile{ExtensionLinearRunSingleShift.lean},
  \leanfile{ExtensionLinearRunSharpness.lean},
  \leanfile{ExtensionLinearRunProofs.lean},
  \leanfile{ExtensionAperiodicSixProofs.lean},
  \leanfile{ExtensionTwoSidedProfileProofs.lean},
  \leanfile{SwapCodeAnswersProofs.lean},
  \leanfile{OperatorDualityProofs.lean},
  \leanfile{NativeSweepInterfaceProofs.lean}, and
  \leanfile{BoundedAritySetCapacityProofs.lean}), proving the two-stage composition theorems and
  finite controls, the one-stage linear run-extension and $\Z_6$ bounds, the two-sided coarse--fine
  and three-term/interval swap-code theorems, the Operator Schmidt Rank bounds and row-shear /
  mutual-information controls, the frozen-fibre cut-causal cone determination theorem, and the
  one-hop bounded-arity set-capacity and clean-entry bounds; and
\item \emph{Seven consequence, transfer, and bridge modules}
  (\leanfile{ExtensionFieldTransfer.lean},
  \leanfile{AffineCubeConsequences.lean},
  \leanfile{SumsetRankConsequences.lean},
  \leanfile{DualityAndSwapConsequences.lean},
  \leanfile{ShuffleShearBridge.lean},
  \leanfile{MachineSweepConsequences.lean}, and
  \leanfile{SetCapacityConsequences.lean}), proving the algebraic-closure primitive-root and
  scalar-extension transfer lemmas, the explicit constants $c(M_1, M_2)$ and
  $D_0(M_1, M_2)$, the general-field finite controls at $d = 13$ and $d = 4$, and the complex
  witness and bridge at $d = 12$, Stone's
  iteration $S^m = \tau$ and the \leanname{TensorProduct} anti-diagonal equivalences, the transport
  of \leanname{LeanMachine.ShuffleShear} to \leanname{Transpose.TM}, the symbolic
  \lean{CausalConeTree} evaluator, the $1/(\Delta+1)$ random-order bound, and the negative
  column-class cut-reduction bounds.
\end{enumerate}

Every one of the $41$ numbered theorems, lemmas, propositions, corollaries, and examples in the
manuscript is completely formalized (all $41$ marked \verb|\formal|) via proved \texttt{theorem}
declarations in these modules. Outside numbered theorem statements, remarks and prose comparisons with
other papers (including the van der Waerden regime comparisons in \Cref{rem:2a-vs-p19} and
\Cref{rem:2b-threshold}(a), the finite crossover and $(u,v)$-minimization comparisons in
\Cref{sec:linear-capacity} checked in exact rational arithmetic in
\leanfile{tools/verify_linear_capacity_crossovers.py}, the general $(m, s)$ padded dimension formula
$4m(m+1) + \frac{s(s+1)}{2} + 2ms$ in \Cref{sec:remark-5-6} verified in Lean at
$(m, d) \in \{(1,4), (1,5), (2,7)\}$, the positive-characteristic Mahler binomial observation in
\Cref{rem:char-dividing-d}, and the numerical search in \Cref{rem:two-sided-numerical}) are proved in
the text (or checked by script) only unless a Lean declaration is named.
